\uuid{tYu8}
\exo7id{164}
\titre{exo7 164}
\auteur{cousquer}
\organisation{exo7}
\datecreate{2003-10-01}
\isIndication{false}
\isCorrection{false}
\chapitre{Logique, ensemble, raisonnement}
\sousChapitre{Récurrence}
\cpge{mpsi-mp2i.raisonnement-vocabulaire-ensembliste.a, pcsi.raisonnement-vocabulaire-ensembliste.a, ptsi.raisonnement-vocabulaire-ensembliste.a, mpsi-mp2i.nombres-reels-suites.i, pcsi.nombres-reels-suites.g, ptsi.nombres-reels-suites.g, bcpst1.logique-ensembles-raisonnement.c, bcpst1.suites-reelles-usuelles}
\module{Algèbre}
\niveau{L1}
\difficulte{}

\contenu{
\texte{
On considère une suite $(u_n)_{n\in \mathbb{N}}$ telle que ~:
$$u_0=0 \quad \mbox{et}\quad u_1=1 \quad\mbox{et}\quad 
\forall n\geq1,\; u_{n+1}=u_n+2u_{n-1}$$
Démontrer que~:
}
\begin{enumerate}
    \item \question{$\forall n\in \mathbb{N},\; u_n\in \mathbb{N}$,}
    \item \question{$\forall n\in \mathbb{N},\; u_n=\frac{1}{3}(2^n-(-1)^n)$.}
\end{enumerate}
}
